\label{ap:codigo_fuente}
\footnotesize 

En este apéndice se exponen los fragmentos de código C\# más significativos desarrollados para el prototipo de \textit{Reset}, ilustrando la implementación de la persistencia de datos, la navegación táctica de la inteligencia artificial, el desacoplamiento de diálogos y las físicas calibradas de salto.

% =========================================================================
\section{Persistencia y Guardado de Datos (SaveManager.cs)}
\label{ap:sec:persistencia}
El siguiente fragmento detalla la primera parte del método para estructurar y serializar el progreso de la partida en formato JSON:

\begin{verbatim}
public void SaveGame()
{
    try
    {
        SaveData data = new SaveData();
        string currentScene = SceneManager.GetActiveScene().name;
        data.escenaGuardada = currentScene;
        data.playerCores = PlayerPrefs.GetInt("PlayerCores", 0);

        // 1. Guardar Inventario Mundo 1
        if (InventarioManager.Instance != null)
        {
            foreach (var slot in InventarioManager.Instance.objetosGuardados)
            {
                if (slot.objeto != null)
                {
                    data.objetosNombres.Add(slot.objeto.nombreObjeto);
                    data.objetosCantidades.Add(slot.cantidad);
                }
            }
        }

        // 2. Guardar Estadísticas de Combate y Salud del Jugador
        PlayerController pc = FindFirstObjectByType<PlayerController>();
        if (pc != null)
        {
            var propArma = pc.GetType().GetField("armaDesbloqueada", 
                System.Reflection.BindingFlags.NonPublic | 
                System.Reflection.BindingFlags.Instance);
            if (propArma != null) 
                data.armaDesbloqueadaMundo1 = (bool)propArma.GetValue(pc);
            data.balasActualesCargadorMundo1 = pc.balasActualesCargador;
        }

        PlayerHealth ph = FindFirstObjectByType<PlayerHealth>();
        if (ph != null)
        {
            var propVida = ph.GetType().GetField("vidaActual", 
                System.Reflection.BindingFlags.NonPublic | 
                System.Reflection.BindingFlags.Instance);
            if (propVida != null) 
                data.vidaActualMundo1 = (int)propVida.GetValue(ph);
        }

        if (currentScene == "01_Hub" || 
            currentScene == "1_Level1" || 
            currentScene == "1_Level2")
        {
            GameObject playerObj = GameObject.FindWithTag("Player");
            if (playerObj != null)
            {
                data.playerPosXMundo1 = playerObj.transform.position.x;
                data.playerPosYMundo1 = playerObj.transform.position.y;
                data.tienePosicionGuardadaMundo1 = true;
            }
        }
\end{verbatim}

\clearpage
El siguiente bloque recopila los identificadores destruidos, las variables estáticas del Mundo 2 y realiza la escritura física del archivo serializado en el disco duro:

\begin{verbatim}
        data.objetosDestruidosMundo1 = new List<string>(destroyedObjects);
        data.dialogosReproducidos = new List<string>(dialogosReproducidos);

        // 3. Guardar estado de Checkpoints y Coleccionables de Mundo 2
        data.livesMundo2 = PlayerPlatformerController.lives;
        data.remainingTimeMundo2 = PlayerPlatformerController.remainingTime;
        data.consecutiveCheckpointDeathsMundo2 = 
            PlayerPlatformerController.consecutiveCheckpointDeaths;
        data.totalCoinsMundo2 = PlayerPlatformerController.totalCoins;
        data.secretCoinsCollectedMundo2 = 
            PlayerPlatformerController.secretCoinsCollected;
        data.collectedCoinsMundo2 = new List<string>(
            PlayerPlatformerController.collectedCoinsActive);
        data.lastCheckpointSceneMundo2 = 
            PlayerPlatformerController.lastCheckpointScene;
        data.lastCheckpointPosXMundo2 = 
            PlayerPlatformerController.lastCheckpointPos.x;
        data.lastCheckpointPosYMundo2 = 
            PlayerPlatformerController.lastCheckpointPos.y;

        PlayerPrefs.SetInt("SavedLevel", 1);
        PlayerPrefs.SetInt("PlayerCores", data.playerCores);
        PlayerPrefs.Save();

        // 4. Escritura JSON física en disco
        string json = JsonUtility.ToJson(data, true);
        File.WriteAllText(saveFilePath, json);
    }
    catch (System.Exception e)
    {
        Debug.LogError("Error al guardar partida: " + e.ToString());
    }
}

public void LoadGame()
{
    if (!HasSaveGame()) return;

    string json = File.ReadAllText(saveFilePath);
    pendingLoadData = JsonUtility.FromJson<SaveData>(json);

    if (SceneTransitionManager.Instance != null)
    {
        SceneTransitionManager.Instance.LoadSceneWithFade(pendingLoadData.escenaGuardada);
    }
    else
    {
        SceneManager.LoadScene(pendingLoadData.escenaGuardada);
    }
}
\end{verbatim}

\clearpage
% =========================================================================
\section{Gestión Modular y Navegación de IA en Mundo 1 (EnemyMelee.cs)}
\label{ap:sec:context_steering}
El fragmento describe la rutina de planificación de estados de persecución y el cálculo del vector de evitación de obstáculos local por direccionamiento de contexto:

\begin{verbatim}
private void EjecutarPersecucion(float distancia)
{
    if (explota && distancia <= rangoActivacionExplosion && !explosionIniciada)
    {
        explosionIniciada = true;
        StartCoroutine(RutinaExplosion());
        return;
    }
    
    if (haceEmbestidas && distancia <= rangoEmbestida && puedeEmbestir && 
        !explosionIniciada)
    {
        if (rb != null) rb.linearVelocity = Vector2.zero;
        StartCoroutine(RutinaEmbestida());
        return;
    }

    if (esZombie && !haceEmbestidas && !explota && 
        distancia <= rangoAtaqueZombie && puedeAtacarZombie && 
        !estaOcupado && !explosionIniciada)
    {
        if (rb != null) rb.linearVelocity = Vector2.zero;
        StartCoroutine(RutinaAtaqueZombie());
        return;
    }
    
    if (!estaOcupado || explosionIniciada) 
    {
        MoverConContextSteering(jugador.position, velocidadNormal, true);
    }
}

private void MoverConContextSteering(Vector2 objetivo, float velocidadBase, bool esPersecucion)
{
    if (rb == null) return;
    Vector2 direccionAlObjetivo = 
        (objetivo - (Vector2)transform.position).normalized;
    
    float[] intereses = new float[8];
    float[] peligros = new float[8];

    // Fase 1: Proyección de Vectores de Interés (Producto Escalar)
    for (int i = 0; i < 8; i++)
    {
        float dot = Vector2.Dot(direcciones[i], direccionAlObjetivo);
        intereses[i] = Mathf.Max(0f, dot);
    }

    // Fase 2: Proyección de Vectores de Peligro (CircleCast)
    for (int i = 0; i < 8; i++)
    {
        Vector2 dir = direcciones[i];
        RaycastHit2D hit = Physics2D.CircleCast(transform.position, radioObstaculos, 
            dir, distanciaRaycast, obstacleLayer);
        float peligro = 0f;

        if (hit.collider != null && hit.collider.gameObject != this.gameObject)
        {
            if (jugador != null && hit.collider.transform == jugador)
                continue;
            peligro = 1f - (hit.distance / distanciaRaycast);
        }
        peligros[i] = peligro;
    }
\end{verbatim}

\clearpage
Lógica de mezcla vectorial, frenesí a corta distancia, pasos rápidos probabilísticos y oscilación sinusoidal lateral:

\begin{verbatim}
    // Fase 3: Fusión Vectorial (Vector Blend)
    Vector2 direccionFinal = Vector2.zero;
    for (int i = 0; i < 8; i++)
    {
        float score = intereses[i] - peligros[i] * factorEvasion;
        score = Mathf.Max(0f, score);
        direccionFinal += direcciones[i] * score;
    }

    if (direccionFinal == Vector2.zero)
        direccionFinal = direccionAlObjetivo;
    else
        direccionFinal = direccionFinal.normalized;
        
    float velocidadActual = velocidadBase;

    if (esZombie)
    {
        if (esPersecucion)
        {
            float distanciaAlJugador = Vector2.Distance(transform.position, jugador.position);
            if (distanciaAlJugador <= distanciaFrenesi)
            {
                velocidadActual *= multiplicadorFrenesi;
            }
            else
            {
                if (tiempoPasoRapidoRestante > 0)
                {
                    tiempoPasoRapidoRestante -= Time.fixedDeltaTime;
                    velocidadActual *= multiplicadorPasoRapido;
                }
                else if (Random.value < probabilidadPasoRapido * Time.fixedDeltaTime)
                {
                    tiempoPasoRapidoRestante = duracionPasoRapido;
                }
            }
        }

        // Modificador errático de balanceo
        Vector2 perpendicular = new Vector2(-direccionFinal.y, direccionFinal.x);
        float factorTambaleo = Mathf.Sin((Time.time + offsetTambaleo) * velocidadTambaleo) 
            * amplitudTambaleo;
        
        direccionFinal = (direccionFinal + perpendicular * factorTambaleo).normalized;
    }

    rb.linearVelocity = direccionFinal * velocidadActual;
}
\end{verbatim}

La corrutina que gobierna la embestida física lineal del enemigo Charger:

\begin{verbatim}
private IEnumerator RutinaEmbestida()
{
    estaOcupado = true;
    puedeEmbestir = false;
    if (rb != null) rb.linearVelocity = Vector2.zero;

    if (spriteRenderer != null) spriteRenderer.color = colorAvisoEmbestida;
    yield return new WaitForSeconds(tiempoPreparacion); 

    if (spriteRenderer != null) spriteRenderer.color = colorOriginal;
    Vector2 posicionObjetivo = jugador.position;
    estaEmbistiendo = true;

    if (anim != null && tieneParametroCharging) anim.SetBool("isCharging", true);

    float maxChargeTime = 2.0f; 
    float elapsed = 0f;
    Vector2 lastPos = transform.position;
    float stuckTimer = 0f;

    while (Vector2.Distance(transform.position, posicionObjetivo) > 0.1f && 
           elapsed < maxChargeTime)
    {
        transform.position = Vector2.MoveTowards(
            transform.position, posicionObjetivo, 
            velocidadEmbestida * Time.deltaTime);
        elapsed += Time.deltaTime;

        if (Vector2.Distance(transform.position, lastPos) < 0.001f)
        {
            stuckTimer += Time.deltaTime;
            if (stuckTimer > 0.15f) break; 
        }
        else
        {
            stuckTimer = 0f;
        }
        lastPos = transform.position;
        yield return null;
    }

    if (anim != null && tieneParametroCharging) anim.SetBool("isCharging", false);
    estaEmbistiendo = false;
    estaOcupado = false;

    yield return new WaitForSeconds(tiempoRecargaEmbestida);
    puedeEmbestir = true;
}
\end{verbatim}

\clearpage
% =========================================================================
\section{Físicas de Plataformas e Inercia Aérea (PlayerPlatformerController.cs)}
\label{ap:sec:fisicas}
El siguiente fragmento detalla la lógica de actualización física de saltos variables y el Coyote Time:

\begin{verbatim}
private void UpdateJumpAndPhysics()
{
    if (isGrounded)
    {
        coyoteTimeCounter = coyoteTime;
        doubleJumpAvailable = true;
    }
    else
    {
        coyoteTimeCounter -= Time.deltaTime;
    }

    bool jumpPressed = Keyboard.current != null && 
        Keyboard.current.spaceKey.wasPressedThisFrame;
    if (jumpPressed)
    {
        jumpBufferCounter = jumpBufferTime;
    }
    else
    {
        jumpBufferCounter -= Time.deltaTime;
    }

    if (coyoteTimeCounter > 0f && jumpBufferCounter > 0f)
    {
        rb.linearVelocity = new Vector2(rb.linearVelocity.x, jumpForce);
        jumpBufferCounter = 0f;
        coyoteTimeCounter = 0f;
    }
\end{verbatim}

\clearpage
El siguiente fragmento detalla el Wall Jump, Double Jump y el frenado vertical para salto variable:

\begin{verbatim}
    else if (enableWallJump && !isGrounded && isTouchingWall && 
             jumpBufferCounter > 0f)
    {
        isWallSliding = false;
        jumpBufferCounter = 0f;
        coyoteTimeCounter = 0f;

        float jumpDir = wallOnRight ? -1f : 1f;
        rb.linearVelocity = new Vector2(
            wallJumpForce.x * jumpDir, wallJumpForce.y);

        if ((jumpDir > 0 && !facingRight) || (jumpDir < 0 && facingRight))
        {
            Flip();
        }
        StartCoroutine(WallJumpControlLockRoutine());
    }
    else if (enableDoubleJump && doubleJumpAvailable && !isGrounded && 
             !isWallSliding && jumpBufferCounter > 0f)
    {
        rb.linearVelocity = new Vector2(rb.linearVelocity.x, doubleJumpForce);
        doubleJumpAvailable = false;
        jumpBufferCounter = 0f;
        coyoteTimeCounter = 0f;
    }

    bool jumpReleased = Keyboard.current != null && 
        Keyboard.current.spaceKey.wasReleasedThisFrame;
    if (jumpReleased && rb.linearVelocity.y > 0f)
    {
        rb.linearVelocity = new Vector2(rb.linearVelocity.x, 
            rb.linearVelocity.y * jumpHeightMultiplier);
        coyoteTimeCounter = 0f;
    }
}
\end{verbatim}

\clearpage
% =========================================================================
\section{Motor de Diálogos Asíncrono en Tiempo de Pausa (DialogueManager.cs)}
\label{ap:sec:dialogos}
Lógica de inicio, avance y encolamiento de frases secuenciales en la interfaz de usuario:

\begin{verbatim}
public void StartDialogue(Dialogue dialogue, System.Action onEnd = null)
{
    if (dialoguePanel == null || nameText == null || 
        dialogueText == null) return;

    Time.timeScale = 0f; 
    dialoguePanel.SetActive(true);
    onDialogueEndCallback = onEnd;

    sentencesQueue.Clear();
    foreach (DialogueLine line in dialogue.lines)
    {
        sentencesQueue.Enqueue(line);
    }
    DisplayNextSentence();
}

public void DisplayNextSentence()
{
    if (isTyping)
    {
        StopCoroutine(typingCoroutine);
        dialogueText.text = currentSentence;
        isTyping = false;
        return;
    }

    if (sentencesQueue.Count == 0)
    {
        EndDialogue();
        return;
    }

    DialogueLine nextLine = sentencesQueue.Dequeue();
    nameText.text = nextLine.npcName;
    currentSentence = nextLine.sentence;
    
    typingCoroutine = StartCoroutine(TypeSentence(currentSentence));
}
\end{verbatim}

\clearpage
Lógica de corrutina en tiempo real y finalización asíncrona de diálogos con callback:

\begin{verbatim}
private IEnumerator TypeSentence(string sentence)
{
    dialogueText.text = "";
    isTyping = true;
    foreach (char letter in sentence.ToCharArray())
    {
        dialogueText.text += letter;
        yield return new WaitForSecondsRealtime(typingSpeed); 
    }
    isTyping = false;
}

public void EndDialogue()
{
    if (dialoguePanel != null) dialoguePanel.SetActive(false);

    Time.timeScale = 1f; 

    if (onDialogueEndCallback != null)
    {
        System.Action callback = onDialogueEndCallback;
        onDialogueEndCallback = null;
        callback.Invoke(); 
    }
}
\end{verbatim}

\clearpage
% =========================================================================
\section{Gestión de Enemigos en Mundo 2 mediante Herencia y Polimorfismo (EnemyPlatformerBase.cs y EnemyPatrolPlatformer.cs)}
\label{ap:sec:enemigos_herencia}

\subsection{Clase Base (EnemyPlatformerBase.cs)}
Lógica base de recepción de daño, animación de muerte controlada y detección de pisada:

\begin{verbatim}
using UnityEngine;

public class EnemyPlatformerBase : MonoBehaviour
{
    public int maxHealth = 1;
    protected int currentHealth;
    public int damageToPlayer = 1;
    [SerializeField] private AudioClip sonidoMuerte;

    protected virtual void Start()
    {
        currentHealth = maxHealth;
    }

    public virtual void TakeDamage(int damage)
    {
        currentHealth -= damage;
        if (currentHealth <= 0) Die();
    }

    protected virtual void Die()
    {
        if (sonidoMuerte != null && AudioManager.Instance != null)
            AudioManager.Instance.PlaySFX(sonidoMuerte);

        Animator anim = GetComponent<Animator>();
        if (anim != null)
        {
            foreach (AnimatorControllerParameter param in anim.parameters)
            {
                if (param.type == AnimatorControllerParameterType.Bool) 
                    anim.SetBool(param.name, false);
                else if (param.type == AnimatorControllerParameterType.Float) 
                    anim.SetFloat(param.name, 0f);
            }
            anim.SetTrigger("Die");
            
            Rigidbody2D rb = GetComponent<Rigidbody2D>();
            if (rb != null) rb.simulated = false; 
            
            this.enabled = false;
            Destroy(gameObject, 0.8f);
        }
        else
        {
            Destroy(gameObject);
        }
    }
\end{verbatim}

\clearpage
Lógica del trigger de colisiones de la clase base:

\begin{verbatim}
    protected virtual void OnCollisionEnter2D(Collision2D collision)
    {
        if (collision.gameObject.GetComponent<EnemyPlatformerBase>() != null)
        {
            TurnAround();
            return;
        }

        if (collision.gameObject.CompareTag("Player"))
        {
            PlayerPlatformerController player = 
                collision.gameObject.GetComponent<PlayerPlatformerController>();
            if (player != null)
            {
                ContactPoint2D contact = collision.GetContact(0);
                if (contact.normal.y <= -0.5f) 
                {
                    player.Bounce();
                    TakeDamage(1);
                }
                else
                {
                    player.TakeDamage(damageToPlayer);
                }
            }
        }
    }

    public virtual void TurnAround() {}
}
\end{verbatim}

\clearpage
\subsection{Clase Hija (EnemyPatrolPlatformer.cs)}
Comportamiento heredado de patrulla con raycasting e inversión de escala local:

\begin{verbatim}
public class EnemyPatrolPlatformer : EnemyPlatformerBase
{
    public float patrolSpeed = 3f;
    public Transform ledgeCheck; 
    public Transform wallCheck;  
    public float checkDistance = 0.5f;
    public LayerMask groundLayer;

    private Rigidbody2D rb;
    private bool movingRight = true;

    protected override void Start()
    {
        base.Start(); 
        rb = GetComponent<Rigidbody2D>();
        movingRight = transform.localScale.x > 0;
    }

    void Update()
    {
        if (patrolSpeed <= 0.01f) return;

        rb.linearVelocity = new Vector2(
            (movingRight ? 1 : -1) * patrolSpeed, rb.linearVelocity.y);

        bool hittingWall = Physics2D.Raycast(wallCheck.position, 
            movingRight ? Vector2.right : Vector2.left, checkDistance, groundLayer);
        bool hittingLedge = Physics2D.Raycast(ledgeCheck.position, 
            Vector2.down, checkDistance, groundLayer);

        if (hittingWall || !hittingLedge)
        {
            TurnAround();
        }
    }

    public override void TurnAround()
    {
        movingRight = !movingRight;
        Vector3 scaler = transform.localScale;
        scaler.x *= -1;
        transform.localScale = scaler;
    }
}
\end{verbatim}